\documentclass[11pt]{article}

\usepackage[T1]{fontenc}
\usepackage[margin=1.1in]{geometry}
\usepackage{amsmath,amssymb,amsthm}
\usepackage{enumitem}
\usepackage{booktabs}
\usepackage[colorlinks=true,linkcolor=blue,citecolor=blue,urlcolor=blue]{hyperref}

\newtheorem{theorem}{Theorem}[section]
\newtheorem{proposition}[theorem]{Proposition}
\newtheorem{lemma}[theorem]{Lemma}
\newtheorem{corollary}[theorem]{Corollary}
\newtheorem{definition}[theorem]{Definition}
\newtheorem{conjecture}[theorem]{Conjecture}
\newtheorem{question}[theorem]{Question}
\theoremstyle{remark}
\newtheorem{remark}[theorem]{Remark}
\newtheorem{example}[theorem]{Example}

\newcommand{\R}{\mathbb{R}}
\newcommand{\X}{X}
\newcommand{\Xhat}{\widehat{\X}}
\newcommand{\Mx}{M_{\X}}
\newcommand{\Rx}{\mathcal{R}_{\X}}
\newcommand{\Hb}[1]{\mathcal{H}_{#1}}
\newcommand{\Sb}[1]{S_{#1}}
\newcommand{\dist}{\operatorname{dist}}
\newcommand{\conv}{\operatorname{conv}}
\newcommand{\hconv}{\operatorname{hconv}}
\newcommand{\Mf}{M_{f}}
\newcommand{\bdry}{\partial_{\infty}}
\newcommand{\sca}[2]{\left\langle #1,\, #2\right\rangle}
\newcommand{\E}{\mathbb{E}}
\newcommand{\HH}{\mathbb{H}}

\title{Medial-axis reconstruction in Hadamard manifolds:\\
a faithful, curvature-free characterization\thanks{Third note in a
series: \cite{R1} generalised the sufficiency half of Theorem~6 of
\cite{Bia} to Hadamard manifolds; \cite{R2} proved the full
characterisation in real hyperbolic space $\HH^{n}$; the Euclidean
case is formalised in \cite{Notes}. Drafted with the assistance of
Claude (Anthropic); all arguments were re-derived and checked by hand,
and every numerical claim verified directly. Remaining errors are
ours.}}
\author{}
\date{June 2026}

\begin{document}
\maketitle

\begin{abstract}
For a closed set $X$ in a complete Riemannian manifold, the
\emph{reconstruction} $\Rx=\bigcup_{a\in\Mx}B(a,r(a))$ is the union of
the open medial balls. Bia\l{}o\.zyt's Theorem~6 characterises $\Rx$
for $X\subseteq\R^{n}$, and a companion note settled real hyperbolic
space; we prove a faithful, unconditional analogue valid in
\emph{every} Hadamard manifold and for \emph{every} closed set:
\[
  \Rx=(\Xhat\setminus X)\ \cup\
  \bigcup\bigl\{\operatorname{int}\Hb{}:\ X\cap\partial\Hb{}\neq
  \hconv(X\cap\partial\Hb{})\bigr\},
\]
where $\Xhat$ is the horoball hull and $\hconv$ is the convex hull cut
out by the ambient \emph{Busemann tilts}. The defect certificate ---
the contact set differs from its horoball hull --- is decidable from
$X$ and the geometry of $M$, with the medial axis absent; it is
$\conv$ on flat horospheres (recovering Theorem~6 verbatim) and the
Carnot hull on homogeneous ones. Both analytic inputs are
curvature-free: a deepest-shadow necessity (Danskin's theorem) and a
swallow lemma using only horofunction convergence. The unconditional
\emph{dynamical} form --- defectiveness read off the medial axis ---
is an immediate consequence of the curvature-free trichotomy; the
theorem above upgrades its certificate to a faithful, $\Mx$-free one.
The remainder of the note makes ``defective'' concrete across a
hierarchy of horospheres governed by the Gauss equation
$K_{S}=K_{M}+\det\mathrm{II}$: flat (warped products $\Mf$ of variable,
pinched curvature --- the first reconstruction theorem beyond constant
curvature), $\mathrm{CAT}(0)$ (warped over a Hadamard fibre), and
positively curved (the Heisenberg horosphere of $\mathbb{CH}^{2}$).
\end{abstract}

\tableofcontents

\bigskip
\noindent\textbf{Part I} develops the main theorem; \textbf{Part II}
makes its defect certificate concrete across the curvature hierarchy.

\section{Introduction}\label{sec:intro}

For a closed set $\X$ in a complete Riemannian manifold $M$ write
$\Mx$ for its \emph{medial axis} (the points with two or more nearest
points of $\X$), $r(a)=\dist(a,\X)$ for the reach of a centre
$a\in\Mx$, and
\[
  \Rx=\bigcup_{a\in\Mx}B\!\left(a,r(a)\right)
\]
for the \emph{reconstruction} of $\X$: the union of the open medial
balls. Bia\l{}o\.zyt's Theorem~6 \cite{Bia} characterises $\Rx$ for
closed $\X\subseteq\R^{n}$: with $\Xhat=\overline{\conv}\,\X$,
\[
  \Rx=(\Xhat\setminus\X)\cup\bigcup\bigl\{\operatorname{int}H:\
  H\ \text{a supporting halfspace},\
  \X\cap\partial H\neq\conv(\X\cap\partial H)\bigr\}.
\]
A companion note \cite{R2} proved the analogue in real hyperbolic
space: for a graph-type closed $\X\subseteq\HH^{n}$,
\begin{equation}\label{eq:hyp-main}
  \Rx=(\Xhat\setminus\X)\ \cup\
  \bigcup\bigl\{\Hb{}\ \text{maximal free, horospherically
  defective}\bigr\},
\end{equation}
where $\Xhat$ is the horoball hull and a maximal free horoball $\Hb{}$
is \emph{horospherically defective} when its contact set
$\Hb{}\cap\X$ is non-convex inside the intrinsically flat horosphere
$\partial\Hb{}\cong\E^{n-1}$.

\subsection*{What was special about hyperbolic space}

The proof of \eqref{eq:hyp-main} reads the geometry of a horoball
$\Hb{}=\{z>1\}$ (upper half-space model) through its flat horosphere
$\{z=1\}\cong\E^{n-1}$. Two facts do all the work.
\begin{enumerate}[label=(\alph*)]
\item For a high centre $a=(\alpha,t)$, the cosine rule
$\cosh d(a,q)=1+\tfrac{1}{2t}\bigl(|\alpha-\kappa|^{2}+(t-1)^{2}\bigr)$
for $q=(\kappa,1)$ is monotone in the \emph{Euclidean} distance
$|\alpha-\kappa|$; so the feet of $a$ are the Euclidean nearest
contacts, and $a$ is medial exactly when its lateral position $\alpha$
has two or more nearest contacts.
\item \textbf{Motzkin's theorem}: a closed set in $\E^{n-1}$ is convex
iff it is Chebyshev. Hence a contact set admits a medial centre iff it
is non-convex.
\end{enumerate}
Both facts use that horospheres of $\HH^{n}$ are \emph{flat}, and the
normalisation ``WLOG $\Hb{}=\{z>1\}$'' uses that
$\operatorname{Isom}(\HH^{n})$ acts transitively on horoballs. These
are precisely the two properties that fail off constant curvature. The
Gauss equation makes the failure quantitative: a horosphere $S$ has
intrinsic curvature $K_{S}=K_{M}+\det(\mathrm{II})$, which vanishes in
$\HH^{n}$ (umbilic, $K_{M}=-1$) but is in general nonzero and can even
be positive under pinching, so horospheres need not be
$\mathrm{CAT}(0)$ and Motzkin can fail.

\subsection*{The main theorem}

The characterization nonetheless survives in full generality, once the
right notion of ``defective'' is found. Reconstruction rests on a
\emph{curvature-free backbone} --- a completeness inclusion, a
rigidity statement, a drift theorem through the horofunction boundary,
and a structure trichotomy --- valid in every Hadamard manifold
(\S\ref{sec:backbone}). The backbone reduces the problem, outside the
hull, to drift into a maximal free horoball; the only remaining task is
to certify \emph{which} horoballs are reached, from $X$ alone. The
correct certificate is not metric. For a maximal free horoball $\Hb{}$
at $\xi$, let $\hconv(K)$ be the convex hull of $K\subseteq\partial\Hb{}$
with respect to the \emph{Busemann tilts} of $M$ --- the trace on
$\partial\Hb{}$ of the intersection of the (tilted) horoballs
containing $K$ (Definition~\ref{def:hconv}). This hull exists in every
Hadamard manifold, needs no horosphere metric and no homogeneity, and
equals $\conv$ on a flat horosphere and the Carnot hull on a
homogeneous one. With it:

\begin{quote}
\textbf{Theorems~\ref{thm:faithful}--\ref{thm:faithful-nc}.} For every
Hadamard manifold $M$ and every closed $\X\subseteq M$,
\[
  \Rx=(\Xhat\setminus X)\ \cup\
  \bigcup\bigl\{\operatorname{int}\Hb{}:\ X\cap\partial\Hb{}\neq
  \hconv(X\cap\partial\Hb{})\bigr\}.
\]
\end{quote}

This is a \emph{faithful} analogue of Theorem~6: its defect
certificate --- the contact set differs from its horoball hull --- is
decidable from $X$ and the geometry of $M$, with the medial axis
$\Mx$ absent from the right-hand side, exactly as
``$X\cap L\neq\conv(X\cap L)$'' is in $\R^{n}$. The two analytic inputs
are curvature-free: necessity is Danskin's first-order condition for
the depth functional --- the deepest free horoball over $p$ has a
non-horoconvex contact set (\S\ref{sec:hconv}) --- and sufficiency is a
\emph{swallow} lemma that uses only horofunction convergence
(Lemma~\ref{lem:swallow}), no curvature bound. Reading defectiveness
\emph{dynamically} --- $\Hb{}$ admits medial centres drifting to $\xi$
--- gives an unconditional characterization for free from the
trichotomy alone (Theorem~\ref{thm:fullgen}); the theorem above is the
upgrade that removes $\Mx$ from the certificate.

\subsection*{Making ``defective'' concrete}

The certificate is intrinsic to $M$, but reading it as a property of
the \emph{horosphere} --- a metric convexity, as in Theorem~6 ---
requires localisation (the feet of a high centre are the
intrinsic-nearest contacts) and confronts the Gauss obstruction.
Part~II carries this out across the curvature hierarchy of horospheres.
The pivot is a family of Hadamard manifolds that is non-homogeneous and
of non-constant, two-sidedly pinched curvature, yet whose horospheres
remain flat: the warped products
$\Mf=(\R\times\R^{n-1},\,dt^{2}+f(t)^{2}g_{\E})$ with $f$ positive,
convex and decreasing to $0$ (\S\ref{sec:warp}). On $\Mf$ the entire
hyperbolic dictionary goes through verbatim --- feet are exactly the
Euclidean nearest contacts (\S\ref{sec:dict}) --- and the
reconstruction dichotomy ``$\Rx\neq\varnothing$ iff the contact set is
non-convex'' holds exactly as in $\HH^{n}$ (\S\ref{sec:recon}), the
warp entering only through a discrimination weight. This is the first
genuinely non-constant-curvature case; the one new effect is that the
up/down symmetry of horoballs breaks (Remark~\ref{rem:asym}). Warping
over a \emph{curved} Hadamard fibre then produces $\mathrm{CAT}(0)$
horospheres, on which the sharpness half of the dichotomy is
unconditional (\S\ref{sec:curved}); and the Heisenberg horosphere of
$\mathbb{CH}^{2}$ gives the first reconstruction theorem on a
\emph{positively} curved horosphere (\S\ref{sec:ch2}). The genuine
frontier --- non-$\mathrm{CAT}(0)$ horospheres in general --- is
isolated in \S\ref{sec:frontier}.

\bigskip
\begin{center}\textsc{Part I. The main theorem}\end{center}

\section{The curvature-free backbone}\label{sec:backbone}

We isolate, once, the part of the theory that uses no curvature
computation beyond $\sec\le0$. Throughout, $M$ is a Hadamard manifold
(complete, simply connected, $\sec\le0$), $\X\subseteq M$ closed, and
$\Xhat$ its horoball hull (the complement of the union of all open
free horoballs).

\begin{theorem}[Reconstruction engine; {\cite{R1,R2}}]
\label{thm:engine}
For every Hadamard manifold $M$ and closed $\X\subseteq M$:
\begin{enumerate}[label=\textup{(\arabic*)}]
\item \emph{(Hull.)} $\Xhat\setminus\X\subseteq\Rx$
  \cite[Thm.~3.1]{R1}.
\item \emph{(Rigidity.)} For $a\in\Mx$ and any foot $q$, no inflation
  $B(a',r')\supsetneq B(a,r(a))$ with $a'$ on the ray from $q$ through
  $a$ is free; medial centres cannot be escaped along a foot-geodesic.
\item \emph{(Drift.)} If $p\in\Rx\setminus\Xhat$ then a sequence of
  medial centres reconstructing $p$ drifts to an ideal point
  $\xi\in\bdry M$ and is captured by the maximal free horoball at
  $\xi$ through $p$, whose horosphere absorbs the feet
  \cite[II.8.13]{BH}.
\item \emph{(Trichotomy.)} Every $p\in M\setminus\X$ is of exactly one
  type: hull, drift, or ray-escape; there is no residual case.
\end{enumerate}
\end{theorem}

Items (1)--(4) are proved for general Hadamard $M$ in \cite{R1,R2} and
quoted here unchanged. The reduction they achieve is: \emph{outside the
hull, reconstruction $\Leftrightarrow$ the drift horoball is reached by
medial centres.} What remains is to turn ``reached by medial centres''
into a certificate read from $X$ alone.

\subsection{The swallow, and the dynamical characterization}

The sufficiency direction --- a reached horoball is reconstructed in
full --- is curvature-free. Call a maximal free horoball $\Hb{}$ with
ideal centre $\xi$ \emph{defective} if $X$ admits medial centres
$a_{k}\in\Mx$ converging to $\xi$ (equivalently $b_{\xi}(a_{k})\to
-\infty$: the centres drift into $\Hb{}$ toward $\xi$); the drift is
\emph{anchored} if the nearest contacts of the $a_{k}$ stay bounded,
and $\Hb{}$ is \emph{anchored defective} if it admits an anchored
drift.

\begin{lemma}[Swallow, curvature-free]\label{lem:swallow}
Let $M$ be \emph{any} Hadamard manifold and $\Hb{}$ an anchored
defective maximal free horoball. Then $\operatorname{int}\Hb{}
\subseteq\Rx$.
\end{lemma}

\begin{proof}
Normalise $b:=b_{\xi}$ so $\Hb{}=\{b<0\}$, $X\subseteq\{b\ge0\}$, and
fix a basepoint $o\in\partial\Hb{}$ (so $b(o)=0$). Drift to $\xi$ is
convergence in the horofunction compactification \cite[II.8.13]{BH}:
for every fixed $y$, $d(y,a_{k})-d(o,a_{k})\to b(y)-b(o)=b(y)$. Let
$x_{k}$ be a nearest contact of $a_{k}$, $r(a_{k})=d(a_{k},x_{k})$; by
anchoring, after a subsequence $x_{k}\to x_{\infty}\in X$.

\emph{Ascent.} For any contact $y$ with $b(y)=\delta$,
$r(a_{k})\le d(a_{k},y)$, so $d(a_{k},x_{k})-d(o,a_{k})\le
d(a_{k},y)-d(o,a_{k})\to\delta$; as $\inf_{X}b=0$ at a maximal free
horoball, $\delta\downarrow0$ gives
$\lim_k[d(a_{k},x_{k})-d(o,a_{k})]\le 0$. Since
$d(a_{k},x_{k})-d(o,a_{k})\to b(x_{\infty})$ (because $x_{k}\to
x_{\infty}$ and $b$ is the limit horofunction) and $b(x_{\infty})\ge0$,
we get $b(x_{\infty})=0$.

\emph{Swallow.} For $p\in\operatorname{int}\Hb{}$, $\tau_{p}:=-b(p)>0$,
\[
  d(p,a_{k})-r(a_{k})=\bigl[d(p,a_{k})-d(o,a_{k})\bigr]
  -\bigl[d(x_{k},a_{k})-d(o,a_{k})\bigr]
  \ \longrightarrow\ b(p)-b(x_{\infty})=-\tau_{p}<0 .
\]
So $d(p,a_{k})<r(a_{k})$ for large $k$: $p\in B(a_{k},r(a_{k}))$ with
$a_{k}\in\Mx$, hence $p\in\Rx$. \emph{No curvature bound is used.}
\end{proof}

\begin{theorem}[Characterization, dynamical form]\label{thm:fullgen}
Let $M$ be \emph{any} Hadamard manifold and $X\subseteq M$ an arbitrary
closed set. With no further hypothesis,
\[
  \Rx\setminus\Xhat\ \subseteq\
  \bigcup\{\operatorname{int}\Hb{}:\Hb{}\ \text{maximal free and
  defective}\}\qquad\text{and}\qquad \Xhat\setminus X\subseteq\Rx ,
\]
and $\operatorname{int}\Hb{}\subseteq\Rx$ for every anchored defective
$\Hb{}$ (Lemma~\ref{lem:swallow}). Consequently, whenever every
defective horoball is anchored or $\sec\le-a^{2}<0$,
\[
  \Rx=(\Xhat\setminus X)\ \cup\
  \bigcup\bigl\{\operatorname{int}\Hb{}:\Hb{}\ \text{maximal free and
  defective}\bigr\} .
\]
This holds in particular for every compact $X$ (any Hadamard $M$), in
the flat case $M=\R^{n}$, and under pinching $\sec\le-a^{2}<0$ --- so
it contains Bia\l{}o\.zyt's Theorem~6 and the pinched theorem as
special cases.
\end{theorem}

\begin{proof}
Each medial ball $B(a,r(a))$ misses $X$, so $\Rx\cap X=\varnothing$;
every open free horoball is disjoint from $\Xhat$, so the right-hand
terms meet neither $X$ nor each other.

\emph{Inclusion ``$\subseteq$''.} Let $p\in\Rx\setminus\Xhat$ (if
instead $p\in\Xhat$ then $p\in\Xhat\setminus X$, as $p\notin X$). By
the trichotomy (Theorem~\ref{thm:engine}(4)) $p$ is hull, drift, or
ray-escape; it is not hull, and not ray-escape (a ray-escape point is
uniquely footed along its escaping ray, hence unreconstructed by
rigidity, Theorem~\ref{thm:engine}(2)). So $p$ is of drift type: its
reconstructing medial centres drift to some $\xi$ with $p$ in the
maximal free horoball $\Hb{}$ at $\xi$ (Theorem~\ref{thm:engine}(3)),
which they witness defective, and $p\in\operatorname{int}\Hb{}$. The
hull inclusion $\Xhat\setminus X\subseteq\Rx$ is completeness
\cite[Thm.~3.1]{R1}.

\emph{Inclusion ``$\supseteq$'' (under the stated condition).} For an
anchored defective $\Hb{}$, $\operatorname{int}\Hb{}\subseteq\Rx$ by
Lemma~\ref{lem:swallow}; under $\sec\le-a^{2}<0$ the same holds for
every defective $\Hb{}$, since horospheres then contract toward $\xi$
(Heintze--Im Hof \cite{HIH}) and the swallow margin survives even when
the nearest contacts escape. Compactness of $X$ forces all contacts,
hence every drift, to be anchored; in $\R^{n}$ a defective halfspace
admits an anchored drift over its Motzkin bisector. Combined with
``$\subseteq$'' this is the boxed equality.
\end{proof}

The dynamical certificate refers to the medial axis, so
Theorem~\ref{thm:fullgen} is an unconditional \emph{skeleton} rather
than a faithful Theorem-6 analogue: it shows $\Rx$ is always the hull
together with a union of defective-horoball interiors, and isolates the
remaining task --- a certificate decidable from $X$. That is the work
of the next section.

\section{The horospherical convex hull, and the main theorem}
\label{sec:hconv}

The faithful certificate replaces ``$\Hb{}$ is reached by medial
centres'' with ``the contact set differs from its hull,'' the hull
taken with respect to the Busemann geometry of $M$. We construct it,
prove its one structural lemma, and obtain the characterization.

\begin{definition}[Horospherical convex hull]\label{def:hconv}
Let $\Hb{}$ be a maximal free horoball at $\xi$, with horosphere
$S=\partial\Hb{}$. For an ideal direction $u$ (a path
$\xi_{\varepsilon}\to\xi$ in $\bdry M$), let
$L_{u}(y)=\frac{d}{d\varepsilon}\big|_{0^{+}}b_{\xi_{\varepsilon}}(y)$
be the one-sided derivative of the Busemann function (it exists, $b$
being convex in $\xi$). The \emph{horospherical convex hull} of
$K\subseteq S$ is
\[
  \hconv(K)=\bigl\{\sigma\in S:\ L_{u}(\sigma)\ge
  \textstyle\inf_{K}L_{u}\ \text{for every }u\bigr\},
\]
the points not separable below $K$ by any Busemann tilt. Equivalently,
$\hconv(K)$ is the trace on $S$ of the \emph{horoball hull} of $K$ ---
the intersection of the tilted horoballs containing $K$, since a tilted
horoball $\{b_{\xi_{\varepsilon}}\ge c\}$ contains $K$ iff
$c\le\inf_{K}b_{\xi_{\varepsilon}}$, and $\sigma$ lies in all of them
iff $L_{u}(\sigma)\ge\inf_{K}L_{u}$ for every $u$. So $\hconv$ is a
genuine \emph{hull}, built from the Busemann geometry of $M$ alone,
with $\Mx$ absent: the certificate $X\cap\partial\Hb{}\neq
\hconv(X\cap\partial\Hb{})$ reads literally as ``the contact set
differs from its hull.'' On a flat horosphere horoballs are
halfspaces, $L_{u}(y)=\langle y,u\rangle$, and $\hconv=\conv$ ---
recovering Bia\l{}o\.zyt's $X\cap L\neq\conv(X\cap L)$ verbatim; on a
homogeneous horosphere it is the Carnot convex hull.
\end{definition}

\begin{lemma}[Tilt functionals are axial invariants]\label{lem:Lu-axial}
For any ideal point $\xi$ and tilt $u$, the functional $L_{u}$ is
constant along every geodesic ray $\gamma$ asymptotic to $\xi$:
$L_{u}(\gamma(t))=L_{u}(\gamma(0))$ for all $t$. In particular, for a
point $p$ over a horosphere foot $\sigma(p)$ on the axis to $\xi$,
$L_{u}(p)=L_{u}(\sigma(p))$.
\end{lemma}

\begin{proof}
The Busemann gradient along $\gamma$ is the inward radial field,
$\nabla b_{\xi}(\gamma(t))=-\dot\gamma(t)$, so
$\frac{d}{dt}b_{\xi_{\varepsilon}}(\gamma(t))=\langle\nabla
b_{\xi_{\varepsilon}}(\gamma(t)),\dot\gamma(t)\rangle$ takes, at
$\varepsilon=0$, the value $-1$ --- the minimum of an inner product of
two unit vectors ($|\nabla b_{\xi_{\varepsilon}}|\equiv1$,
$|\dot\gamma|=1$). An inner extremum has vanishing
$\varepsilon$-derivative, so, exchanging the smooth mixed derivatives,
\[
  \frac{d}{dt}L_{u}(\gamma(t))
  =\frac{d}{d\varepsilon}\Big|_{0}\frac{d}{dt}
  b_{\xi_{\varepsilon}}(\gamma(t))=0,
\]
and $L_{u}\circ\gamma$ is constant. In $\R^{n}$,
$L_{u}(y)=\langle y,u^{\perp}\rangle$ with $u^{\perp}\perp\xi$, visibly
constant along the $\xi$-axis.
\end{proof}

\begin{theorem}[Faithful, unconditional --- compact $X$]
\label{thm:faithful}
Let $M$ be \emph{any} Hadamard manifold and $X\subseteq M$ compact.
Then
\[
  \boxed{\ \Rx=(\Xhat\setminus X)\ \cup\
  \bigcup\bigl\{\operatorname{int}\Hb{}:\ X\cap\partial\Hb{}\neq
  \hconv(X\cap\partial\Hb{})\bigr\}\ },
\]
the defect certificate --- the contact set is not horospherically
convex --- being decidable from $X$ and the geometry of $M$, with
$\Mx$ absent. It reduces to Bia\l{}o\.zyt's
$X\cap\partial\Hb{}\neq\conv(\cdot)$ on flat horospheres ($\R^{n}$,
$\HH^{n}$) and to its Carnot form on homogeneous ones.
\end{theorem}

\begin{proof}
\emph{Faithfulness.} Every object on the right is computed from $X$
and the horoballs of $M$, never from the medial axis: the hull
$\Xhat$ is the intersection of the free horoballs, the maximal free
horoballs $\Hb{}$ are determined by $\{\Hb{}\cap X=\varnothing\}$, the
contact set is $X\cap\partial\Hb{}$, and $\hconv$ is the horoball hull
of that contact set (Definition~\ref{def:hconv}). So the certificate
is decidable by looking at $X$ and its hull, with $\Mx$ absent.

\emph{``$\subseteq$''.} The hull term is $\Xhat\setminus X\subseteq\Rx$
\cite[Thm.~3.1]{R1}. Let $p\in\Rx\setminus\Xhat$ and let $\xi^{\ast}$
maximise the depth $D(\xi)=\inf_{x\in X}b_{\xi}(x)-b_{\xi}(p)$
(attained, as $X$ is compact), with deepest free horoball
$\Hb{}^{\ast}$, contacts $K^{\ast}$, and foot $\sigma(p)$ of $p$ on
$\partial\Hb{}^{\ast}$. The infimum $\inf_{x}b_{\xi}(x)$ is attained on
$K^{\ast}$, so by Danskin's theorem the one-sided derivative of $D$
along a tilt $u$ is $\inf_{K^{\ast}}L_{u}-L_{u}(\sigma(p))$
(Lemma~\ref{lem:Lu-axial} gives $L_{u}(p)=L_{u}(\sigma(p))$, as $p$ and
$\sigma(p)$ lie on one ray to $\xi^{\ast}$); maximality forces it
$\le0$ for every $u$, i.e.\ $\sigma(p)\in\hconv(K^{\ast})$. And
$\sigma(p)\notin K^{\ast}$: were $\sigma(p)$ a contact $x_{0}$, every
$x\in X$ would have $d(p,x)\ge b_{\xi^{\ast}}(x)-b_{\xi^{\ast}}(p)\ge
D(\xi^{\ast})=d(p,x_{0})$, with equality only along the axis, so
$x_{0}$ is the unique foot of $p$ and the distance-flow carries $p$ up
the axis to infinity, contradicting $p\in\Rx$ \cite{Lie}. Hence
$K^{\ast}\neq\hconv(K^{\ast})$ and $p\in\operatorname{int}\Hb{}^{\ast}$.

\emph{``$\supseteq$''.} Let $\Hb{}$ be maximal free with
$K:=X\cap\partial\Hb{}\neq\hconv(K)$, and pick
$\sigma\in\hconv(K)\setminus K$. A centre high over $\sigma$ lies in
$\operatorname{int}\Hb{}\subseteq M\setminus\Xhat$, so it is not of
hull type; nor of ray-escape type, for $\hconv$, taken over all tilts
$\pm u$, surrounds $\sigma$ on every side: escape toward $u$ would need
every contact behind $\sigma$ ($L_{u}<L_{u}(\sigma)$), whereas the
$-u$ constraint $L_{-u}(\sigma)\ge\inf_{K}L_{-u}$ gives a contact
\emph{ahead} ($L_{u}\ge L_{u}(\sigma)$) blocking it. By the trichotomy
(Theorem~\ref{thm:engine}(4)) the centre is of drift type, so $\Hb{}$
admits drifting medial centres and is defective; the swallow
(Lemma~\ref{lem:swallow}, anchored as $X$ is compact, or along the
medial ray) gives $\operatorname{int}\Hb{}\subseteq\Rx$. \emph{No
Motzkin or metric convexity in the horosphere is invoked: the
sufficiency is the curvature-free trichotomy, and $\hconv$ enters only
as the hull that rules out escape.}
\end{proof}

The two analytic inputs are both curvature-free: the deepest-shadow is
Danskin's first-order condition for the depth functional (it is
\cite{R2}'s lemma for $\HH^{n}$, here recognised as convex
optimisation valid in any Hadamard manifold, with $\conv$ replaced by
$\hconv$), and the swallow is Lemma~\ref{lem:swallow}. The remaining
inputs are orthogonal to curvature: the compactness of $X$ --- removed
next --- and the regularity of $\xi^{\ast}$ on the ideal boundary,
where $L_{u}$ is a one-sided derivative.

\begin{theorem}[Faithful, unconditional --- arbitrary closed $X$]
\label{thm:faithful-nc}
Theorem~\ref{thm:faithful} holds for \emph{every} closed $X\subseteq M$,
with $\hconv$ read on the extended contact set
$\overline{K}=K\cup K_{\infty}$, where $K_{\infty}$ collects the ideal
points of the horosphere $\partial\Hb{}$ that $X$ approaches
asymptotically along it. ($\hconv(\overline{K})$ adds to $\hconv(K)$
exactly the recession directions: $\inf_{\overline K}L_{u}=-\infty$
wherever $X$ recedes, voiding the constraint, as in Bia\l{}o\.zyt's
recession cone.) The inclusion ``$\subseteq$'' and the defect
certificate are unconditional; the equality holds in full whenever the
radial offset of the drift vanishes --- under any pinching, in the
flat case, and along a bisector ray, hence for every model in this
note. The single exception is a purely asymptotic drift
($K=\varnothing$) in an \emph{inhomogeneous} $\sec\le0$ manifold where
the offset does not vanish: there only the deep part
$\{\tau_{p}>\limsup_{k}\mathrm{off}_{k}\}$ of $\operatorname{int}\Hb{}$
is reconstructed by this drift, the defect itself still detected.
\end{theorem}

\begin{proof}
\emph{Necessity, with no compactness.} The depth
$D(\xi)=\inf_{x\in X}b_{\xi}(x)-b_{\xi}(p)$ is upper semicontinuous on
the compact ideal boundary $\bdry M$ --- an infimum of the continuous
maps $\xi\mapsto b_{\xi}(x)$ minus the continuous $b_{\xi}(p)$ --- so it
attains its maximum at some $\xi^{\ast}$ for any closed $X$, removing
the single use of compactness in Theorem~\ref{thm:faithful}. At
$\xi^{\ast}$ the Danskin derivative of $D$ along $u$ is
$\bigl(\inf_{x\in X^{\ast}}L_{u}(x)\bigr)-L_{u}(\sigma(p))$, where
$X^{\ast}$ is the set of finite or ideal minimisers of
$b_{\xi^{\ast}}$ over $\overline{X}$; maximality forces it $\le0$ for
all $u$, giving $\sigma(p)\in\hconv(\overline{K}^{\ast})$, and
screening (a finite foot $\sigma(p)\in K^{\ast}$ is the unique nearest
point of $p$, ideal contacts being infinitely far, so $p$ escapes up
the axis) gives $\sigma(p)\notin K^{\ast}$. Hence $K^{\ast}\neq
\hconv(\overline{K}^{\ast})$.

\emph{Sufficiency.} Let $\sigma\in\hconv(\overline{K})\setminus K$. A
centre high over $\sigma$ lies in
$\operatorname{int}\Hb{}\subseteq M\setminus\Xhat$, so it is not of
hull type; nor of ray-escape type, for $\hconv(\overline{K})$, taken
over all tilts $\pm u$, surrounds $\sigma$ on every side: in any
candidate escape direction $\eta$ there is an extended contact ahead
--- a finite one, or an ideal one in $K_{\infty}$ along which $X$
approaches $\partial\Hb{}$, so the receding ray re-encounters $X$ and
the distance to $X$ cannot grow without bound. By the trichotomy
(Theorem~\ref{thm:engine}(4)) the centre is of drift type, so $\Hb{}$
admits medial centres $a_{k}\to\xi$ and is defective.

It remains to reconstruct $\operatorname{int}\Hb{}$ from such $a_{k}$,
whose feet may escape to $K_{\infty}$. With $\tau_{k}:=-b(a_{k})$ and a
basepoint $o\in\partial\Hb{}$: $X\subseteq\{b\ge0\}$ and $1$-Lipschitz
$b$ give $r(a_{k})\ge\tau_{k}$, while horofunction convergence gives,
for $p\in\operatorname{int}\Hb{}$,
$d(p,a_{k})=\tau_{k}+\mathrm{off}_{k}-\tau_{p}+o(1)$, where
$\mathrm{off}_{k}=d(o,a_{k})+b(a_{k})\ge0$ is the radial offset of
$a_{k}$. Thus $r(a_{k})-d(p,a_{k})\ge\tau_{p}-\mathrm{off}_{k}-o(1)$.
When $\sec\le-a^{2}<0$ (asymptotic rays converge) or $M$ is flat
(Pythagoras) --- hence on $\R^{n}$, $\HH^{n}$, every $\Mf$, and every
pinched manifold --- $\mathrm{off}_{k}\to0$, so $r-d\to\tau_{p}>0$ and
$\operatorname{int}\Hb{}\subseteq\Rx$ in full. Equivalently, where the
horosphere supports a Motzkin (or Carnot) bisector witness one may take
the $a_{k}$ along the ray over it, $\mathrm{off}_{k}\equiv0$ and feet
escaping; the valley $\{t=-1/(1+x^{2})\}$ in $\Mf$ ($K=\varnothing$,
feet $x^{\ast}\to\pm\infty$) is so reconstructed.
\end{proof}

\begin{remark}[The offset, resolved in the final note]
\label{rem:offset-resolved}
The final note of the series \cite{C} removes the single exception by
changing the normalisation. Defining a \emph{medial limit function} as
a locally uniform limit of the free-ball functions
$d(\cdot,a_{k})-r(a_{k})$ over medial $a_{k}$ --- normalised by their
own reach, rather than against a Busemann function chosen in advance
--- the swallow becomes unconditional: $\{\psi<0\}\subseteq\Rx$ for
every such limit $\psi$, with no anchoring, pinching, or offset
hypothesis \cite[Lem.~4.1]{C}. A drift of asymptotic offset $c$ is
then exactly a medial limit at \emph{level} $c$ ($\psi=b_{\xi}+c$),
reconstructing $\{b_{\xi}<-c\}$ in full; the deep-part phenomenon
above is thereby the precise content of the level, not a gap in the
reconstruction. What remains open is only the faithful computation of
the minimal level $c(\xi)$ from $X$ in an inhomogeneous manifold
\cite[Rem.~4.6]{C}; it vanishes in every case named in the theorem.
\end{remark}

\subsection{The static reading: localisation and Motzkin}
\label{sec:static}

Theorems~\ref{thm:faithful}--\ref{thm:faithful-nc} are complete and
faithful, but their certificate is phrased through the ambient
$\hconv$. To read it as Theorem~6 does --- a convexity of the contact
set in the \emph{intrinsic} geometry of the horosphere --- one needs to
identify $\hconv$ with a metric hull, which requires \emph{localisation}.

\begin{definition}[Defective, static]\label{def:defective}
A maximal free horoball $\Hb{}$ is \emph{defective} if its contact set
$X\cap\partial\Hb{}$ is \emph{not Chebyshev} in the intrinsic geometry
of the horosphere $\partial\Hb{}$: some point of $\partial\Hb{}$ has
two or more nearest points of the contact set. When $\partial\Hb{}$ is
flat, Motzkin's theorem (Theorem~\ref{thm:motzkin}) makes this
equivalent to the contact set being non-convex; on a curved
horosphere, non-Chebyshev is the correct invariant, in general
strictly weaker than non-convex.
\end{definition}

Two comparison facts connect this static notion to the dynamical and
ambient ones:
\begin{itemize}
\item[\textup{(L)}] \emph{Localisation:} among contacts on
$\partial\Hb{}$, those nearest to a high centre are those nearest in
the intrinsic metric of $\partial\Hb{}$ --- so a centre is medial iff
its position has two or more intrinsic-nearest contacts;
\item[\textup{(NC)}] \emph{No collapse:} the two feet of a drifting
medial centre converge to \emph{distinct} contacts.
\end{itemize}
Under (L) the static $\hconv$ is the metric convex hull and
``defective'' (Definition~\ref{def:defective}) agrees with the
ambient certificate of Theorem~\ref{thm:faithful}; (NC), needed only
for the converse, is automatic where geodesics diverge (flat or
$\mathrm{CAT}(0)$ horospheres) and can fail where they focus
(positively curved horospheres). \emph{Motzkin is not part of the
theorem:} the invariant is non-Chebyshev, read off the two feet;
Motzkin only restates it as non-convex when the horosphere is flat.

The obstruction to (L), and to a metric reading at all, is the
intrinsic curvature of the horosphere. Part~II quantifies it (the Gauss
equation) and then establishes the static certificate concretely on
flat, $\mathrm{CAT}(0)$, and positively curved horospheres.

\bigskip
\begin{center}\textsc{Part II. ``Defective'' across the curvature
hierarchy}\end{center}

\section{The Gauss obstruction}\label{sec:gauss}

\begin{proposition}[Horosphere curvature]\label{prop:gauss}
Let $S$ be a horosphere in a Hadamard manifold $M$, with (outward)
shape operator $\mathrm{II}\succeq0$. For a tangent $2$-plane
$\Pi\subseteq T S$ spanned by an orthonormal pair $u,v$,
\[
  K_{S}(\Pi)=K_{M}(\Pi)
  +\bigl(\mathrm{II}(u,u)\,\mathrm{II}(v,v)-\mathrm{II}(u,v)^{2}\bigr).
\]
In $\HH^{n}$ every horosphere is umbilic with $\mathrm{II}=\mathrm{Id}$
and $K_{M}\equiv-1$, so $K_{S}\equiv0$. Under
$-b^{2}\le\sec\le-a^{2}<0$ the principal curvatures of every horosphere
lie in $[a,b]$, hence $K_{S}\in[a^{2}-b^{2},\,b^{2}-a^{2}]$.
\end{proposition}

\begin{proof}
The displayed identity is the Gauss equation for the hypersurface
$S\subseteq M$. Horospheres are limits of geodesic spheres, whose shape
operators converge to that of the stable Jacobi tensor; the bound
$\mathrm{II}\in[a,b]\cdot\mathrm{Id}$ (as quadratic forms) under
$-b^{2}\le\sec\le-a^{2}$ is the Heintze--Im Hof estimate \cite{HIH}.
Combining $K_{M}(\Pi)\in[-b^{2},-a^{2}]$ with
$\det(\mathrm{II}|_{\Pi})\in[a^{2},b^{2}]$ gives the stated range.
\end{proof}

The point is the sign change: when $b>\sqrt{2}\,a$ the upper end
$b^{2}-a^{2}$ exceeds $a^{2}$ and the interval contains positive
values, so a horosphere can have positive intrinsic curvature in some
$2$-planes. On a space of positive curvature, Chebyshev sets need not
be convex, and the Motzkin equivalence has no reason to hold.
\emph{This, not the loss of homogeneity, is the essential obstacle.}
The natural first step is to switch it off: keep the horosphere flat
while letting the ambient curvature vary. The slices $\{t=c\}$ of the
warped products below are umbilic with shape operator
$(f'/f)\,\mathrm{Id}$, and the Gauss equation reads
$K_{S}=-(f'/f)^{2}+(f'/f)^{2}=0$ for every admissible $f$; constant
ambient curvature is the strictly stronger requirement
$(\log f)''=0$. The warped family is thus exactly the room between
``flat horospheres'' and ``constant curvature.''

\section{Warped products with flat horospheres}\label{sec:warp}

Fix $n\ge2$ and a \emph{warp} $f:\R\to(0,\infty)$ that is smooth,
strictly convex, strictly decreasing, with $f(t)\to0$ as
$t\to+\infty$. (Convex, positive and decreasing to $0$ forces
$f(t)\to+\infty$ as $t\to-\infty$.) Let
$\Mf=(\R\times\R^{n-1},\,dt^{2}+f(t)^{2}g_{\E})$, with
$g_{\E}=\sum_{i}(dx^{i})^{2}$.

\begin{proposition}[$\Mf$ is a pinched Hadamard manifold]
\label{prop:mf-curv}
$\Mf$ is a complete, simply connected Riemannian manifold. Its
sectional curvatures are
\[
  K(\partial_{t},\partial_{x^{i}})=-\frac{f''}{f}<0,
  \qquad
  K(\partial_{x^{i}},\partial_{x^{j}})=-\Bigl(\frac{f'}{f}\Bigr)^{2}<0
  \quad(i\neq j),
\]
and every sectional curvature lies in the convex hull of these two
values; in particular $\sec<0$, so $\Mf$ is Hadamard. If
$a^{2}\le f''/f\le b^{2}$ and $a^{2}\le(f'/f)^{2}\le b^{2}$ on all of
$\R$, then $\Mf$ is pinched: $-b^{2}\le\sec\le-a^{2}$. The curvature is
non-constant unless $f(t)=Ce^{kt}$, i.e.\ unless $\Mf\cong\HH^{n}$.
\end{proposition}

\begin{proof}
$\R\times\R^{n-1}=\R^{n}$ is simply connected, and $g_{f}$ is complete
because $f$ is bounded below by a positive constant on every compact
$t$-interval and the $t$-lines have infinite length in both
directions. The curvature formulas are the standard warped-product
identities for a one-dimensional base $B=\R$ and flat fibre
$F=\R^{n-1}$ \cite[Prop.~7.42]{One}:
\[
  K(\partial_{t},Y)=-\frac{f''}{f},\qquad
  K(Y,Z)=\frac{K^{F}-|f'|^{2}}{f^{2}}=-\frac{(f')^{2}}{f^{2}}
  \quad(Y,Z\ \text{fibre, }K^{F}=0).
\]
Because the base is one-dimensional and the fibre flat, the curvature
operator is diagonal in the splitting
$T\Mf=\R\partial_{t}\oplus(\text{fibre})$, so every sectional curvature
is a convex combination of the two displayed eigenvalues; both are
negative under the standing hypotheses. The pinching is immediate. The
curvature is constant iff $f''/f=(f'/f)^{2}$, i.e.\ $(\log f)''=0$,
i.e.\ $f=Ce^{kt}$; then $g_{f}$ is the half-space metric of $\HH^{n}$
after rescaling.
\end{proof}

\begin{example}\label{ex:warp}
$f(t)=e^{-t}(1+\varepsilon e^{-t})$ with $\varepsilon>0$ is admissible.
Here $-f''/f=-\dfrac{1+4\varepsilon e^{-t}}{1+\varepsilon e^{-t}}$
ranges over $(-4,-1)$ and $-(f'/f)^{2}=-\Bigl(\dfrac{1+2\varepsilon
e^{-t}}{1+\varepsilon e^{-t}}\Bigr)^{2}$ over $(-4,-1)$: the manifold
is pinched $-4\le\sec\le-1$, genuinely non-constant, asymptotically
hyperbolic of curvature $-1$ as $t\to+\infty$ and $-4$ as
$t\to-\infty$.
\end{example}

\begin{proposition}[Flat horospheres]\label{prop:mf-horo}
Let $\xi^{+}\in\bdry\Mf$ be the ideal endpoint of the upward
$t$-lines. The Busemann function of $\xi^{+}$ is $b_{\xi^{+}}=-t$ (up
to an additive constant); consequently the horospheres centred at
$\xi^{+}$ are exactly the slices $\Sb{c}=\{t=c\}$. Each $\Sb{c}$
carries the induced metric $f(c)^{2}g_{\E}$ and is therefore
intrinsically flat, isometric to $\E^{n-1}$. The slices are \emph{not}
horospheres for any other ideal point: the Busemann function toward the
downward direction depends on the fibre coordinate (its level sets are
curved), so $\{t=c\}$ is a flat horosphere for $\xi^{+}$ alone.
\end{proposition}

\begin{proof}
Let $\gamma(s)=(s,0)$ be the upward unit-speed $t$-line. For
$y=(t_{y},x_{y})$, $b_{\xi^{+}}(y)=\lim_{s\to\infty}(d(y,\gamma(s))-s)$.
By Lemma~\ref{lem:distlaw} below (whose proof is independent of this
proposition),
$d(y,\gamma(s))=(s-t_{y})+\dfrac{|x_{y}|^{2}}{2I(t_{y},s)}+o(1)$, and
$I(t_{y},s)=\int_{t_{y}}^{s}f^{-2}\to\infty$ as $s\to\infty$ because
$f(t)\to0$ makes $f^{-2}$ non-integrable at $+\infty$. Hence the middle
term vanishes and $b_{\xi^{+}}(y)=-t_{y}$, independent of $x_{y}$. Its
level sets are the slices $\{t=c\}$, with induced metric
$g_{f}|_{t=c}=f(c)^{2}g_{\E}$, a constant multiple of the flat metric.
Toward the downward axis $\gamma(-s)=(-s,0)$ the same computation gives
$d(y,(-s,0))-s\to t_{y}+\dfrac{|x_{y}|^{2}}{2\int_{-\infty}^{t_{y}}
f^{-2}}$, where now $\int_{-\infty}^{t_{y}}f^{-2}$ is \emph{finite}
(since $f\to\infty$, $f^{-2}$ is integrable at $-\infty$). The
surviving $|x_{y}|^{2}$-term shows the downward Busemann level sets are
paraboloidal, not the slices; the slices are horospheres for $\xi^{+}$
only.
\end{proof}

\begin{proposition}[Transitivity on horospheres]\label{prop:mf-isom}
For every $v\in\R^{n-1}$ and every $R\in O(n-1)$, the map
$(t,x)\mapsto(t,Rx+v)$ is an isometry of $\Mf$ fixing each horosphere
$\Sb{c}$ setwise and fixing $\xi^{\pm}$. The group of such maps acts
transitively on each $\Sb{c}$ and by the full group of Euclidean
isometries on its intrinsic geometry. Reflection of $\R^{n-1}$ across
any affine hyperplane is likewise an isometry; its fixed-point set is a
totally geodesic copy of $\Mf$ of one dimension lower.
\end{proposition}

\begin{proof}
$g_{f}=dt^{2}+f(t)^{2}g_{\E}$ is invariant under any map acting as the
identity on $t$ and as a $g_{\E}$-isometry on the fibre, because
$f(t)^{2}g_{\E}$ is then preserved pointwise. Euclidean isometries of
$\R^{n-1}$ are exactly such maps and act transitively on each fibre.
The fixed-point set of the reflection across a hyperplane
$H\subseteq\R^{n-1}$ is $\R\times H$, the fixed-point set of an
isometry, hence totally geodesic; with the restricted metric it is
$M_{f}$ of dimension $\dim H+1$.
\end{proof}

Homogeneity is thus retained \emph{along} horospheres --- the part the
dictionary needs (``a centre may sit above any prescribed lateral
position'') --- while it is lost \emph{across} them: no isometry moves
$\Sb{0}$ to $\Sb{1}$, since $f(0)\neq f(1)$ scales the intrinsic
diameter. This is the precise sense in which $\Mf$ relaxes hyperbolic
homogeneity without touching the flatness of horospheres.

\section{The local dictionary on \texorpdfstring{$\Mf$}{Mf}}
\label{sec:dict}

The two dictionary facts of $\HH^{n}$ now reappear, the first
\emph{exactly} and the second through Motzkin on the flat horosphere.

\begin{lemma}[Distance law]\label{lem:distlaw}
For $a=(t_{a},\alpha)$ and $q=(t_{0},\kappa)$ in $\Mf$ with
$t_{a}>t_{0}$, the distance $d(a,q)$ depends only on $t_{a},t_{0}$ and
$\rho:=|\alpha-\kappa|$; the geodesic between them stays in the totally
geodesic vertical plane through $\alpha,\kappa$; and $d(a,q)$ is a
strictly increasing function of $\rho$ at every scale. Writing
$I=I(t_{0},t_{a})=\int_{t_{0}}^{t_{a}}f(t)^{-2}\,dt$, one has the
rigorous two-sided bound
\[
  (t_{a}-t_{0})\ \le\ d(a,q)\ \le\ (t_{a}-t_{0})+\frac{\rho^{2}}{2I}.
\]
In particular, for fixed $\rho$ and $t_{a}\to+\infty$ (so
$I\to\infty$), $d(a,q)=(t_{a}-t_{0})+\dfrac{\rho^{2}}{2I}(1+o(1))
=(t_{a}-t_{0})+o(1)$.
\end{lemma}

\begin{proof}
Dependence on $\rho$ alone is Proposition~\ref{prop:mf-isom}: a fibre
isometry carries $(\alpha,\kappa)$ to any pair with the same gap.
Reflection across the hyperplane bisecting $\alpha,\kappa$ fixes the
vertical plane $P=\R\times(\kappa+\R(\alpha-\kappa))$ pointwise, so $P$
is totally geodesic and contains the minimiser; the problem is
two-dimensional, $P\cong(\R^{2},dt^{2}+f^{2}dx^{2})$. Strict
monotonicity in $\rho$ holds at all scales by first variation: the
variation field of the endpoint $\alpha$ along the fibre is
$\partial_{x}/|\partial_{x}|$, and $\partial_{\rho}d=\cos\theta>0$,
$\theta<\tfrac{\pi}{2}$ being the angle the incoming geodesic makes
with the outward fibre direction (it approaches from the $\kappa$-side,
by the reflection symmetry).

Parametrise a $t$-monotone path in $P$ by $t\mapsto(t,x(t))$ with
$\int_{t_{0}}^{t_{a}}x'=\rho$, of length
$L[x]=\int_{t_{0}}^{t_{a}}\sqrt{1+f(t)^{2}x'(t)^{2}}\,dt$. The
integrand is $\ge1$, so $L\ge t_{a}-t_{0}$: the lower bound. For the
upper bound use the comparison path of constant momentum
$f^{2}x'\equiv\rho/I$, which satisfies
$\int x'=\frac{\rho}{I}\int f^{-2}=\rho$; by
$\sqrt{1+s}\le1+\tfrac{s}{2}$,
\[
  d(a,q)\le L=\int_{t_{0}}^{t_{a}}\sqrt{1+\frac{\rho^{2}}{f^{2}I^{2}}}
  \,dt\le(t_{a}-t_{0})+\frac{\rho^{2}}{2I^{2}}\int_{t_{0}}^{t_{a}}f^{-2}
  =(t_{a}-t_{0})+\frac{\rho^{2}}{2I}.
\]
For the refined asymptotic, Cauchy--Schwarz on \emph{any} path gives
$\rho^{2}=(\int x')^{2}\le I\int f^{2}x'^{2}$, so along the true
minimiser $\int f^{2}x'^{2}\ge\rho^{2}/I$; with
$\sqrt{1+s}\ge1+\tfrac{s}{2}-\tfrac{s^{2}}{8}$ this yields
$d-(t_{a}-t_{0})\ge\tfrac{\rho^{2}}{2I}-\tfrac18\int(f^{2}x'^{2})^{2}$.
The quartic is negligible: on the minimiser $f^{2}x'\equiv\rho/I$ is
constant, so $f^{2}x'^{2}=\rho^{2}/(f^{2}I^{2})\le\rho^{2}/I$ (since
$f(t)^{-2}\le I$ on $[t_{0},t_{a}]$, $f$ being decreasing), whence
$\int(f^{2}x'^{2})^{2}\le(\rho^{2}/I)\int f^{2}x'^{2}\le(\rho^{2}/I)^{2}
=o(\rho^{2}/I)$ as $\rho^{2}/I\to0$. Combined with the upper bound this
is the stated expansion.
\end{proof}

\begin{lemma}[Feet are exactly the Euclidean nearest contacts]
\label{lem:feet}
Let $\X\subseteq\Mf$ and $a=(t_{a},\alpha)$ with the nearest points of
$\X$ lying on a horosphere $\Sb{0}=\{t=0\}$ --- e.g.\ $\X=K\subseteq
\Sb{0}$. Then the feet of $a$ are exactly the points $(0,\kappa)$,
$\kappa\in K$, minimising the Euclidean distance $|\alpha-\kappa|$.
Consequently $a\in\Mx$ iff $\alpha$ has two or more Euclidean nearest
points in $K$, and the medial directions of $K$ are precisely its
non-Chebyshev points.
\end{lemma}

\begin{proof}
By Lemma~\ref{lem:distlaw}, $d(a,(0,\kappa))$ is a strictly increasing
function of $|\alpha-\kappa|$ for fixed $t_{a},t_{0}=0$; hence its
minimisers over $\kappa\in K$ are exactly the minimisers of
$|\alpha-\kappa|$. Multivaluedness of the nearest-point map is
therefore multivaluedness of Euclidean nearest point in $K$, i.e.\
$\alpha$ non-Chebyshev for $K$.
\end{proof}

This is the entire dictionary. The flat-horosphere hypothesis enters
\emph{only} through Lemma~\ref{lem:feet}; once feet are Euclidean
nearest contacts, Motzkin supplies the equivalence non-convex
$\Leftrightarrow$ admits a medial centre, with no further reference to
the ambient curvature.

\begin{theorem}[Motzkin {\cite{Mot}}]\label{thm:motzkin}
A nonempty closed set $K\subseteq\E^{N}$ is convex if and only if it is
Chebyshev (every point has a unique nearest point of $K$).
\end{theorem}

\section{The reconstruction theorem on \texorpdfstring{$\Mf$}{Mf}}
\label{sec:recon}

\begin{theorem}[Non-convex contact set reconstructs the thin-end
horoball]\label{thm:mf-suff}
Let $K\subseteq\Sb{0}=\{t=0\}$ be closed and non-convex in
$\Sb{0}\cong\E^{n-1}$, and let $\X\supseteq K$ be any closed set with
$\X\cap\{t>-\nu\}=K$ for some $\nu>0$ (in particular $\X=K$). Then the
entire thin-end horoball interior $\{t>0\}$ lies in $\Rx$.
\end{theorem}

\begin{proof}
Since $K$ is non-convex it is non-Chebyshev (Motzkin), so there is a
point $\alpha^{\ast}\in\R^{n-1}$ with two Euclidean nearest points
$k_{1},k_{2}\in K$ at common distance
$\rho:=|\alpha^{\ast}-k_{1}|=|\alpha^{\ast}-k_{2}|$, and no point of
$K$ nearer. For each height $t_{a}>0$ put
$a_{t_{a}}=(t_{a},\alpha^{\ast})$. By Lemma~\ref{lem:feet} the feet of
$a_{t_{a}}$ are $(0,k_{1}),(0,k_{2})$, both at ambient distance
$r(a_{t_{a}})=d(a_{t_{a}},(0,k_{1}))$, so $a_{t_{a}}\in\Mx$ for every
$t_{a}>0$. (Contacts of $\X$ with $t\le-\nu$ are farther: they add at
least $\nu$ to the $t$-gap by Lemma~\ref{lem:distlaw}, for $t_{a}$
large.)

Fix an interior point $p=(t_{p},\pi)$ with $t_{p}>0$. The integrals
$I(0,t_{a})$ and $I(t_{p},t_{a})$ both tend to $+\infty$ as
$t_{a}\to+\infty$ (non-integrability of $f^{-2}$ at $+\infty$), so the
sandwich of Lemma~\ref{lem:distlaw} gives
\[
  r(a_{t_{a}})=d(a_{t_{a}},(0,k_{1}))=t_{a}+o(1),
  \qquad
  d(p,a_{t_{a}})\le(t_{a}-t_{p})+\frac{|\pi-\alpha^{\ast}|^{2}}
  {2I(t_{p},t_{a})}=(t_{a}-t_{p})+o(1).
\]
Hence $d(p,a_{t_{a}})-r(a_{t_{a}})\le-t_{p}+o(1)\to-t_{p}<0$, so for
all large $t_{a}$, $p\in B(a_{t_{a}},r(a_{t_{a}}))$ with
$a_{t_{a}}\in\Mx$, i.e.\ $p\in\Rx$. As $p$ was an arbitrary interior
point, $\{t>0\}\subseteq\Rx$.
\end{proof}

\begin{theorem}[Reconstruction of a set on a horosphere]
\label{thm:mf-sharp}
Let $\X=K\subseteq\Sb{0}=\{t=0\}$ be closed in $\Mf$.
\begin{enumerate}[label=\textup{(\alph*)}]
\item If $K$ is convex in $\Sb{0}\cong\E^{n-1}$, then $\Mx=\varnothing$
  and $\Rx=\varnothing$.
\item If $K$ is non-convex, then $\{t>0\}\subseteq\Rx$.
\end{enumerate}
In particular $\Rx\neq\varnothing\iff K$ is non-convex, the same
dichotomy as in $\HH^{n}$ \textup{\cite[``$\X$ on a horosphere'']{R2}}.
\end{theorem}

\begin{proof}
(b) is Theorem~\ref{thm:mf-suff}. For (a), suppose $K$ is convex, hence
Chebyshev in $\E^{n-1}$ (Motzkin). Let $a=(t_{a},\alpha)$ be arbitrary.
By Lemma~\ref{lem:distlaw}, $d(a,(0,\kappa))$ is a strictly increasing
function of $|\alpha-\kappa|$ at every height $t_{a}$ --- above, on, or
below $\Sb{0}$ --- and at every scale, the monotonicity there being
unconditional. Therefore the nearest point of $a$ in $K$ is the unique
Euclidean nearest point of $\alpha$ in $K$, which exists and is unique
by the Chebyshev property. Thus every centre has a unique foot,
$\Mx=\varnothing$, and $\Rx=\varnothing$.
\end{proof}

\begin{remark}[Up/down asymmetry]\label{rem:asym}
Because $f$ is not even in $t$, the manifold is not $t$-reversible. By
Proposition~\ref{prop:mf-horo} the slices are horospheres for
$\xi^{+}$ \emph{only}; the slab $\{t<0\}$ is the exterior of the
horoball $\{t\ge0\}$, not a horoball. So a contact set
$K\subseteq\Sb{0}$ bounds a genuine flat-horosphere horoball only on
the upward side $\{t>0\}$, which when defective is reconstructed in
full (Theorem~\ref{thm:mf-suff}), the upward weight $1/I(0,t_{a})\to0$.
Toward the downward direction the same $K$ sits on a \emph{curved}
horosphere (the downward Busemann level sets are paraboloidal), and its
reconstruction there is governed by that curved horosphere, not the
flat slice. Concretely, for $K=\{\pm1\}\subseteq\R\subseteq\Sb{0}$ in
Example~\ref{ex:warp}, the upward horoball $\{t>0\}$ is reconstructed
in full, while the slab $\{t<0\}$ is not (off-axis points at depth
escape every axial medial ball) --- as it must be, since the genuine
downward horoball is the curved one. That a single set presents a flat
horosphere to one ideal point and a curved horosphere to another is the
local face of the Gauss obstruction, and it motivates the curved theory
of \S\ref{sec:curved}--\S\ref{sec:frontier}.
\end{remark}

\begin{theorem}[Necessity on $\Mf$, unconditional]\label{thm:mf-nec}
Let $\X\subseteq\Mf$ be closed and $p\in\Rx\setminus\Xhat$. Then $p$
lies in a maximal free horoball whose contact set is non-convex in its
flat horosphere (or, toward a downward ideal point, fails to be
horospherically convex).
\end{theorem}

\begin{proof}
This is the necessity half (``$\subseteq$'') of
Theorem~\ref{thm:faithful}, with Theorem~\ref{thm:faithful-nc} for
non-compact $\X$, specialised to $\Mf$ --- a curvature-free argument
requiring no warped-metric Busemann expansion: the depth
$D(\xi)=\inf_{x\in\X}b_{\xi}(x)-b_{\xi}(p)$ attains its maximum at a
deepest $\xi^{\ast}$, Danskin's theorem with Lemma~\ref{lem:Lu-axial}
gives $\sigma(p)\in\hconv(K^{\ast})$, and screening gives
$\sigma(p)\notin K^{\ast}$, so $K^{\ast}\neq\hconv(K^{\ast})$. When
$\xi^{\ast}=\xi^{+}$ the horosphere is flat
(Proposition~\ref{prop:mf-horo}), the tilt functionals are Euclidean,
and $\hconv=\conv$, so $K^{\ast}\neq\conv(K^{\ast})$: the deepest
horoball is non-convex in its flat horosphere.
\end{proof}

\section{Curved horospheres I: warping over a Hadamard fibre}
\label{sec:curved}

The family $\Mf$ kept horospheres flat by warping over a flat fibre.
Warping over a \emph{curved} fibre produces genuinely curved
horospheres while preserving the exact dictionary, as long as the
curvature stays nonpositive --- so the horospheres are
$\mathrm{CAT}(0)$. This is the first curved-horosphere case, and it
locates the next obstruction precisely: the Motzkin equivalence splits
into an \emph{easy} half (convex $\Rightarrow$ Chebyshev), which
$\mathrm{CAT}(0)$ supplies for free, and a \emph{hard} half (Chebyshev
$\Rightarrow$ convex), the Riemannian Chebyshev-set problem.

Let $N$ be a Hadamard manifold of dimension $n-1$ and $f$ as before.
Set $M_{f,N}=(\R\times N,\,dt^{2}+f(t)^{2}g_{N})$, which reduces to
$\Mf$ when $N=\E^{n-1}$.

\begin{proposition}[Geometry of $M_{f,N}$]\label{prop:fN}
$M_{f,N}$ is a Hadamard manifold. Its horospheres centred at $\xi^{+}$
are the slices $\{t=c\}$, with induced metric $f(c)^{2}g_{N}$: a
rescaled copy of $N$, hence $\mathrm{CAT}(0)$, of sectional curvature
$K_{N}/f(c)^{2}$. The ambient sectional curvatures are $-f''/f$ (planes
through $\partial_{t}$) and $\bigl(K_{N}-(f')^{2}\bigr)/f^{2}\le0$
(planes tangent to the fibre).
\end{proposition}

\begin{proof}
The curvature formulas are O'Neill's for the warped product
$\R\times_{f}N$ \cite[Prop.~7.42]{One}; $K_{N}\le0$ and $(f')^{2}\ge0$
make the fibre curvature $\le0$, and $f''>0$ makes the radial curvature
$<0$. $\R\times N\cong\R^{n}$ is simply connected and complete with
$\sec\le0$, hence Hadamard. The Busemann computation of
Proposition~\ref{prop:mf-horo} used only $f\to0$ and gives
$b_{\xi^{+}}=-t$ and slices $\{t=c\}$ as horospheres here too; their
induced metric is $f(c)^{2}g_{N}$, a homothety of $N$, with curvature
scaled by $f(c)^{-2}$.
\end{proof}

\begin{lemma}[Dictionary over a reflective fibre]\label{lem:fN-dict}
Suppose every geodesic of $N$ is the fixed-point set of an isometric
involution of $N$ (call $N$ \emph{geodesically reflective}; this holds
for $N=\R^{k}\times\HH^{m_{1}}\times\cdots\times\HH^{m_{r}}$, since
geodesics of a Riemannian product are products of geodesics and each
factor reflects across a geodesic). For $a=(t_{a},\alpha)$ and
$q=(t_{0},\kappa)$ the minimising geodesic lies in the totally geodesic
slab $\R\times\gamma$, where $\gamma$ is the $N$-geodesic through
$\alpha,\kappa$; consequently
\[
  (t_{a}-t_{0})\ \le\ d(a,q)\ \le\
  (t_{a}-t_{0})+\frac{d_{N}(\alpha,\kappa)^{2}}{2\,I(t_{0},t_{a})},
\]
and $d(a,q)$ is a strictly increasing function of the intrinsic
distance $d_{N}(\alpha,\kappa)$ at every scale and height. Hence the
feet of a centre are exactly the $d_{N}$-nearest contacts, and a centre
is medial iff its fibre position has two or more nearest contacts in
the intrinsic metric of $N$.
\end{lemma}

\begin{proof}
The reflection $r_{\gamma}$ of $N$ across $\gamma$ is an isometry
fixing $\gamma$ pointwise; $\mathrm{id}_{\R}\times r_{\gamma}$ is then
an isometry of $M_{f,N}$ whose fixed-point set is $\R\times\gamma$,
hence totally geodesic and containing the minimiser. On
$\R\times\gamma$ the metric is $dt^{2}+f(t)^{2}\,ds^{2}$ with $s$ the
$N$-arclength along $\gamma$ --- exactly the two-dimensional warped
metric of Lemma~\ref{lem:distlaw} with horizontal gap
$d_{N}(\alpha,\kappa)$ --- so the bounds and strict monotonicity
transfer verbatim.
\end{proof}

\begin{theorem}[Reconstruction over a Hadamard fibre]\label{thm:fN}
Let $N$ be geodesically reflective as in Lemma~\ref{lem:fN-dict} (e.g.\
$N=\R^{k}\times\HH^{m}$) and $K\subseteq S_{0}=N$ closed.
\begin{enumerate}[label=\textup{(\alph*)}]
\item \emph{(Sharpness, unconditional.)} If $K$ is convex in $N$, then
$\Mx=\varnothing$ and $\Rx=\varnothing$.
\item \emph{(Sufficiency.)} If $K$ is non-Chebyshev in $N$ --- some
point has two or more nearest points --- then $\{t>0\}\subseteq\Rx$.
\end{enumerate}
Hence $\Rx\neq\varnothing\iff K$ is non-Chebyshev in $N$. If $N$ has
the \emph{Motzkin property} (every closed non-convex set is
non-Chebyshev), this reads $\Rx\neq\varnothing\iff K$ non-convex ---
the hyperbolic dichotomy, now on curved $\mathrm{CAT}(0)$ horospheres.
\end{theorem}

\begin{proof}
(a) A convex set in the Hadamard (hence $\mathrm{CAT}(0)$) manifold $N$
is Chebyshev: nearest-point projection onto a complete convex subset of
a $\mathrm{CAT}(0)$ space is single-valued \cite[II.2.4]{BH}. By
Lemma~\ref{lem:fN-dict} the foot of \emph{any} centre
$a=(t_{a},\alpha)$ is the unique $d_{N}$-nearest contact to $\alpha$
(the monotonicity holds at every height), so $\Mx=\varnothing$ and
$\Rx=\varnothing$. (b) A point $\alpha^{\ast}$ with two $d_{N}$-nearest
contacts $k_{1},k_{2}$ gives, by Lemma~\ref{lem:fN-dict}, a two-footed
centre $(t_{a},\alpha^{\ast})$ for every $t_{a}>0$; the swallow
estimate of Theorem~\ref{thm:mf-suff} is purely radial (it uses only
$I(t_{p},t_{a}),I(0,t_{a})\to\infty$) and transfers, giving
$\{t>0\}\subseteq\Rx$.
\end{proof}

Two ingredients carry the sharpness proof: the foot-localisation of
Lemma~\ref{lem:fN-dict} (proved under geodesic reflectivity, and
expected asymptotically for any Hadamard fibre), and the
single-valuedness of projection onto a convex set --- the only place
the \emph{curvature} of the horosphere enters, and precisely the
$\mathrm{CAT}(0)$ property. So, granting localisation, the sharpness
half of the dichotomy survives into curved horospheres exactly as far
as $\mathrm{CAT}(0)$ reaches. The remaining content --- sufficiency
from non-convexity --- is the Riemannian Chebyshev-set problem in $N$,
Motzkin's theorem in $\E^{n-1}$ and open in full generality. The
reconstruction theorem on $M_{f,N}$ is thus \emph{equivalent} to the
Motzkin property of the fibre $N$, isolating the entire difficulty in a
single classical-type question about $N$ alone. Necessity for general
$\X$ is again the curvature-free deepest-shadow of
Theorem~\ref{thm:faithful} ($M_{f,N}$ is Hadamard), needing no separate
shadow lemma; the sharpness statement (a) is unconditional.

\section{Curved horospheres II: complex hyperbolic space}
\label{sec:ch2}

The first \emph{positively} curved horosphere is the Heisenberg
horosphere of $\mathbb{CH}^{2}$. Because the complex-hyperbolic
distance is explicit, localisation --- fact (L) of \S\ref{sec:static},
elsewhere only expected --- can here be carried out \emph{exactly}, and
it yields the first reconstruction theorem on a non-$\mathrm{CAT}(0)$
horosphere.

\paragraph{Horospherical coordinates.} Realise $\mathbb{CH}^{2}$ as the
negative lines in $(\mathbb{C}^{3},\langle\cdot,\cdot\rangle)$ for the
second Hermitian form $\langle z,w\rangle=z_{1}\bar w_{3}+z_{2}\bar
w_{2}+z_{3}\bar w_{1}$ of signature $(2,1)$. A point with Heisenberg
coordinates $(\zeta,v)\in\mathbb{C}\times\R$ and height $u>0$ lifts to
\[
  \psi(\zeta,v,u)=\Bigl(\tfrac{-|\zeta|^{2}-u+iv}{2},\ \zeta,\ 1\Bigr),
  \qquad \langle\psi,\psi\rangle=-u .
\]
The horospheres centred at the distinguished ideal point
($u\to\infty$) are the slices $\{u=c\}$, each a copy of the Heisenberg
group $\mathfrak{H}_{3}=\mathbb{C}\times\R$; the horoball is $\{u>c\}$.

\begin{proposition}[Exact distance, and localisation to the Carnot
metric]\label{prop:ch2}
For $p=(\zeta,v,u)$ and $p'=(\zeta',v',u')$,
\[
  \cosh^{2}\!\tfrac{d(p,p')}{2}
  =\frac{1}{uu'}\!\left[\Bigl(\tfrac{|\zeta-\zeta'|^{2}+u+u'}{2}\Bigr)^{2}
  +\Bigl(\tfrac{v-v'}{2}+\operatorname{Im}(\zeta\bar\zeta')\Bigr)^{2}\right].
\]
Writing $D=|\zeta_{a}-\zeta_{q}|^{2}$ and
$V=v_{a}-v_{q}+2\operatorname{Im}(\zeta_{a}\bar\zeta_{q})$, a high
centre $a=(\zeta_{a},v_{a},u_{a})$ and a contact $q$ on $\{u=u_{0}\}$
obey the \emph{exact} identity
\[
  \cosh^{2}\!\tfrac{d(a,q)}{2}
  =\frac{u_{a}}{4u_{0}}+\frac{D+u_{0}}{2u_{0}}
  +\frac{(D+u_{0})^{2}+V^{2}}{4u_{a}u_{0}} .
\]
Consequently the feet of a high centre are the contacts of least
\emph{horizontal} distance $|\zeta_{a}-\zeta_{q}|$, ties broken by
least $|V|$: the vertical coordinate is suppressed by a factor
$1/u_{a}$. The horizontal/vertical crossover is at
$V\sim\sqrt{2u_{a}D}$, so under the parabolic dilation
$(\zeta,v)\mapsto(\zeta,v/\sqrt{u_{a}})$ the nearest-point structure
converges to that of the Carnot--Carath\'eodory metric on
$\mathfrak{H}_{3}$ --- \emph{not} the Riemannian induced metric
$4|d\zeta|^{2}/u_{0}+\theta^{2}/u_{0}^{2}$, which would weight $V$ at
order one.
\end{proposition}

\begin{proof}
With $\psi_{3}=\psi'_{3}=1$, $\langle\psi,\psi'\rangle=\psi_{1}+\bar
\psi'_{1}+\zeta\bar\zeta' =-\tfrac{|\zeta-\zeta'|^{2}+u+u'}{2}
+i\bigl(\tfrac{v-v'}{2}+\operatorname{Im}(\zeta\bar\zeta')\bigr)$, and
$\cosh^{2}(d/2)=\langle\psi,\psi'\rangle\langle\psi',\psi\rangle/
(\langle\psi,\psi\rangle\langle\psi',\psi'\rangle)
=|\langle\psi,\psi'\rangle|^{2}/(uu')$, which is the first display. The
second is its expansion in $u_{a}$ at fixed $u_{0}$ (no approximation:
$(D+u_{a}+u_{0})^{2}+V^{2}$ over $4u_{a}u_{0}$). The contact dependence
is $\tfrac{D}{2u_{0}}$ at order $1$ and $\tfrac{V^{2}}{4u_{a}u_{0}}$ at
order $u_{a}^{-1}$; equating them gives the crossover $V^{2}=2u_{a}D$.
The induced horosphere metric follows from the infinitesimal form of
the distance: $ds^{2}=4|d\zeta|^{2}/u_{0}+\theta^{2}/u_{0}^{2}$ with
$\theta=dv-2\operatorname{Im}(\zeta\,\overline{d\zeta})$. (Numerically,
as $u_{a}$ runs $3\to1000$ the ambient-nearest of fixed contacts
migrates from the Riemannian-nearest to the horizontal-nearest, and the
crossover ratio $V^{\ast}/\sqrt{2u_{a}D}\to1$.)
\end{proof}

\begin{theorem}[Reconstruction in $\mathbb{CH}^{2}$: sufficiency]
\label{thm:ch2}
Let $K$ be a closed contact set on a horosphere $\{u=u_{0}\}$ of
$\mathbb{CH}^{2}$ that is \emph{graded non-Chebyshev}: some position
admits two contacts of equal least horizontal distance and equal $V$.
Then the whole horoball interior $\{u>u_{0}\}$ lies in $\Rx$.
\end{theorem}

\begin{proof}
Normalise $u_{0}=1$. Graded non-Chebyshev gives
$(\zeta^{\ast},v^{\ast})$ and two contacts $q_{1},q_{2}$ with equal
$D=\delta^{2}$ and equal $V$, hence (Proposition~\ref{prop:ch2})
equidistant from $a_{t}=(\zeta^{\ast},v^{\ast},t)$ for every $t$, with
no contact nearer; so $a_{t}$ is medial for all large $t$. For an
interior point $p=(\zeta_{p},v_{p},u_{p})$, $u_{p}>1$, the distance
formula gives $\cosh^{2}(d(a_{t},q_{1})/2)=\tfrac{t}{4}+O(1)$ and
$\cosh^{2}(d(p,a_{t})/2)=\tfrac{t}{4u_{p}}+O(1)$, whence, via
$\operatorname{arccosh}(x)=\log 2x+o(1)$,
\[
  r(a_{t})=d(a_{t},q_{1})=\log t+O(\tfrac1t),\qquad
  d(p,a_{t})=\log t-\log u_{p}+o(1),
\]
so $d(p,a_{t})-r(a_{t})\to-\log u_{p}<0$. Thus $p\in B(a_{t},r(a_{t}))$
with $a_{t}$ medial for all large $t$, i.e.\ $p\in\Rx$. As $p$ was
arbitrary, $\{u>1\}\subseteq\Rx$.
\end{proof}

Proposition~\ref{prop:ch2} settles localisation in $\mathbb{CH}^{2}$
exactly, with the intrinsic metric the Carnot one, and
Theorem~\ref{thm:ch2} is, to our knowledge, the first medial-axis
reconstruction on a positively curved horosphere. The detection
invariant is graded non-Chebyshev, read off the two feet, so it needs
no Chebyshev-in-Carnot theory. What remains for the full
$\mathbb{CH}^{2}$ characterization: necessity (a drifting medial centre
forces graded non-Chebyshev); the exact reconstructed region for a
general contact set; and, only for restating ``graded non-Chebyshev''
as a convexity, the right notion of convexity in $\mathfrak{H}_{3}$
(horizontal versus Euclidean).

\section{The frontier: non-$\mathrm{CAT}(0)$ horospheres}
\label{sec:frontier}

The dividing line is now sharp. Over a flat fibre ($\Mf$) the
dictionary is literally Motzkin; over a nonpositively curved fibre
($M_{f,N}$, $\mathrm{CAT}(0)$ horospheres) the sharpness half is free
and the dichotomy reduces to the Chebyshev property of the fibre. The
first place where even the sharpness half can fail is a
\emph{positively} curved horosphere --- $K_{S}>0$ in some plane, which
Proposition~\ref{prop:gauss} permits under pinching and which no warped
product over a Hadamard fibre can produce.

\paragraph{Rank-one symmetric spaces.} In $\mathbb{CH}^{n}$ (and
$\mathbb{HH}^{n},\mathbb{OH}^{2}$, Damek--Ricci spaces) the isometry
group still acts transitively on horoballs, so ``WLOG one horoball''
survives. But the horosphere is the Heisenberg group
$\mathfrak{H}_{2n-1}$ with a left-invariant metric \cite{HIH,One,Gol}:
the geodesic up the axis contracts the $2n-2$ horizontal (contact)
directions like $e^{-\tau}$ and the one vertical (centre) direction
like $e^{-2\tau}$ --- the anisotropy behind the $-4\le\sec\le-1$
pinching of $\mathbb{CH}^{n}$. Its left-invariant metric has sectional
curvatures of \emph{both} signs (for $\mathfrak{H}_{3}$, Milnor's
$-\tfrac34$ in the horizontal plane and $+\tfrac14$ in planes through
the centre \cite{Mil}), so it is not $\mathrm{CAT}(0)$. For $n=2$ both
dictionary steps are carried out in \S\ref{sec:ch2} --- localisation
exactly, sufficiency from graded non-Chebyshev. In general $n$,
localisation should give the same with the anisotropic two-rate
Heisenberg metric (the analogue of Lemma~\ref{lem:fN-dict} with a
two-rate contraction in place of the single $1/I$), and detection
becomes a Chebyshev-set problem in a Carnot group, sensitive to which
convexity --- Euclidean, geodesic, or horizontal --- one means.

\paragraph{Pinched negative curvature.} Dropping homogeneity, under
$-b^{2}\le\sec\le-a^{2}$ each horoball must be treated in place, its
horosphere an Alexandrov space with two-sided curvature bounds
(Proposition~\ref{prop:gauss}) and $K_{S}$ generically of both signs.
Localisation becomes a comparison statement --- the discrimination
weight lies between the Heintze--Im Hof rates $e^{-b\tau}$ and
$e^{-a\tau}$, from Jacobi-field control of $\mathrm{II}$ --- and
detection is Motzkin's theorem in an Alexandrov space of indefinite
curvature, the precise admissible class being the open problem.

\paragraph{The stratification.} The phenomenon is therefore controlled,
at every level, by a single intrinsic question about the horosphere ---
whether a contact set with a two-footed centre is non-convex, and
conversely --- and the curvature of the horosphere governs how much of
that equivalence survives:
\begin{itemize}
\item \emph{flat} ($\Mf$): Motzkin, both halves --- full theorem
  (\S\ref{sec:recon});
\item \emph{$\mathrm{CAT}(0)$} ($M_{f,N}$): sharpness free, hard half
  $=$ Motzkin in $N$ --- sharpness theorem (\S\ref{sec:curved});
\item \emph{non-$\mathrm{CAT}(0)$} ($\mathbb{CH}^{2}$): localisation
  and sufficiency via graded non-Chebyshev in the Carnot metric
  (\S\ref{sec:ch2});
\item \emph{non-$\mathrm{CAT}(0)$, general} ($\mathbb{CH}^{n}$, $n\ge3$;
  pinched): two-rate / Alexandrov detection --- open.
\end{itemize}
Flat and $\mathrm{CAT}(0)$ horospheres keep the sharpness half
unconditionally; positive horosphere curvature is the genuine frontier.

\begin{remark}[Prospects for the clean static form]\label{rem:prospects}
The faithful characterization (Theorems~\ref{thm:faithful}--%
\ref{thm:faithful-nc}) is already unconditional; what the frontier
governs is only its \emph{static} reading
(Definition~\ref{def:defective}) --- ``defective $=$ contact set
non-Chebyshev in the intrinsic metric.'' Localisation (L) is expected
throughout the pinched range by Jacobi-field comparison, with one
subtlety: the intrinsic metric may be sub-Riemannian (Carnot) rather
than Riemannian --- a refinement of the \emph{definition} of defective,
not an obstruction, since the invariant is non-Chebyshev, which needs
only a metric. The genuine nub is no-collapse (NC): on a positively
curved horosphere the squared intrinsic distance loses $2$-convexity
past conjugate points, so two feet distinct at every finite height may
coalesce in the limit. The cure, indicated by the hyperbolic
experience, is to read defectiveness at \emph{finite height}: two
distinct feet at each $t$ is all that medial means, and collapse in the
limit is irrelevant, at the price of a per-height rather than
closed-form criterion. Whether the closed-form static criterion itself
survives in full positive-curvature generality reduces to one question:
\emph{can a limit-collapsing medial centre reconstruct?} The decisive
test case is $\mathbb{CH}^{2}$, whose single homogeneous Heisenberg
horosphere makes the computation explicit.
\end{remark}

\paragraph{Conclusion.} A characterization holds in full generality
(Theorem~\ref{thm:fullgen}): for every Hadamard manifold and every
closed set, $\Rx$ is the hull together with the interiors of the
defective maximal free horoballs, defectiveness read dynamically and
the proof resting only on the curvature-free trichotomy and a
curvature-free swallow lemma. Its faithful, $\Mx$-free upgrade
(Theorems~\ref{thm:faithful}--\ref{thm:faithful-nc}) replaces the
dynamical certificate by ``the contact set differs from its horoball
hull $\hconv$,'' which is $\conv$ on flat horospheres (Theorem~6
verbatim) and unconditional in every Hadamard manifold, necessity being
Danskin's deepest-shadow and sufficiency the trichotomy. What is
special to constant curvature is not the theorem but the \emph{static}
reading of ``defective''; the medial-axis reconstruction phenomenon is
\emph{not} special to constant curvature, with the curvature of the
horosphere --- not of the ambient manifold --- governing how concrete
that static reading can be made. On the warped family $\Mf$ the entire
hyperbolic dictionary and the reconstruction dichotomy transfer
verbatim despite variable pinched ambient curvature; warping over a
curved Hadamard fibre gives $\mathrm{CAT}(0)$ horospheres on which the
sharpness half remains an unconditional theorem; and the Heisenberg
horosphere of $\mathbb{CH}^{2}$ gives the first reconstruction theorem
on a positively curved horosphere. The genuine frontier is detection on
non-$\mathrm{CAT}(0)$ horospheres in general: a Chebyshev question in a
Carnot group, or Motzkin's theorem in an Alexandrov space of indefinite
curvature.

\begin{thebibliography}{9}

\bibitem{Bia} A.~Bia\l{}o\.zyt, \emph{Sets reconstructible with medial
axis}, preprint, AGH University of Krak\'ow, 2024.

\bibitem{R1} \emph{Sets reconstructible with medial axis: a
generalisation to Riemannian manifolds}, companion analysis notes,
June 2026 (file \texttt{medial\_axis\_riemannian\_1.pdf}).

\bibitem{R2} \emph{Medial-axis reconstruction in hyperbolic space,
with a structural framework for Hadamard manifolds}, companion note,
June 2026 (file \texttt{medial\_axis\_riemannian\_2.pdf}).

\bibitem{Notes} \emph{Formalization notes on Theorem~6 of
\cite{Bia}}, June 2026 (files
\texttt{medial\_formalization\_notes.tex},
\texttt{medial\_formalization\_plan.md}); Lean~4 formalization on
branch \texttt{medial-axis}, files
\texttt{Mathlib/Analysis/Convex/MedialAxis*.lean}.

\bibitem{C} \emph{Medial-axis reconstruction on complete Riemannian
manifolds: capture, drift, and the unconditional theorem}, companion
note, July 2026 (file \texttt{medial\_axis\_complete.pdf}).

\bibitem{BH} M.~Bridson and A.~Haefliger, \emph{Metric Spaces of
Non-Positive Curvature}, Grundlehren 319, Springer, 1999.

\bibitem{One} B.~O'Neill, \emph{Semi-Riemannian Geometry with
Applications to Relativity}, Pure and Applied Mathematics 103,
Academic Press, 1983. (Warped-product curvature, Prop.~7.42.)

\bibitem{HIH} E.~Heintze and H.-C.~Im Hof, \emph{Geometry of
horospheres}, J.\ Differential Geom.\ \textbf{12} (1977), 481--491.

\bibitem{Mil} J.~Milnor, \emph{Curvatures of left invariant metrics on
Lie groups}, Advances in Math.\ \textbf{21} (1976), 293--329.

\bibitem{Gol} W.~M.~Goldman, \emph{Complex Hyperbolic Geometry}, Oxford
Mathematical Monographs, Oxford University Press, 1999.

\bibitem{Mot} T.~Motzkin, \emph{Sur quelques propri\'et\'es
caract\'eristiques des ensembles born\'es non convexes}, Atti Accad.\
Naz.\ Lincei \textbf{21} (1935), 773--779.

\bibitem{Lie} A.~Lieutier, \emph{Any open bounded subset of
$\mathbb{R}^{n}$ has the same homotopy type as its medial axis},
Computer-Aided Design \textbf{36} (2004), 1029--1046.

\end{thebibliography}

\end{document}
